%% frag_related.tex
%% Subsections to ADD to Section 2 (Related Work) of tosem_amara.tex.
%% Suggested insertion order:
%%   - "Query Rewriting and Expansion" goes immediately AFTER the existing
%%     \subsection{Retrieval-Augmented Generation}.
%%   - "Rewriting for Recall, Reranking for Precision" follows it.
%%   - "Multi-Agent RAG and Reproducibility Infrastructure" goes AFTER the
%%     existing \subsection{Multi-Agent RAG Frameworks} (it extends, and does
%%     not duplicate, that subsection).
%%   - "Evaluation Methodology for Retrieval-Augmented Systems" and
%%     "Negative and Mixed Results" go LAST in Section 2, after
%%     \subsection{Self-Improving Agents}.
%% New \bibitem entries are in frag_related_bibitems.tex.


\subsection{Query Rewriting and Expansion}

The standard justification for rewriting a query before retrieval is vocabulary
mismatch: the terms a person uses for an information need frequently differ from
those in the documents that would satisfy it, an effect documented long before
retrieval became neural~\cite{furnas1987vocabulary}. Classical information
retrieval addressed it by expanding the query with terms drawn from an initial
retrieved set --- pseudo-relevance feedback in the Rocchio
tradition~\cite{rocchio1971relevance} --- or from corpus-wide co-occurrence
statistics~\cite{xu1996localglobal}, a design space surveyed by Carpineto and
Romano~\cite{carpineto2012survey}.

Large language models made the expansion text cheap to generate rather than
mine. HyDE prompts a model for a hypothetical answer document and retrieves
using that document's embedding in place of the query's~\cite{gao2023hyde}.
Query2doc generates a pseudo-document and \emph{concatenates} it with the
original query, retaining the user's terms in the issued
query~\cite{wang2023query2doc}; Jagerman et al.\ compare prompting strategies
for term expansion against pseudo-relevance
feedback~\cite{jagerman2023expansion}. Within RAG,
Rewrite-Retrieve-Read trains a small rewriter to adapt the query to a frozen
retriever and reader~\cite{ma2023rewrite}, DMQR-RAG issues several rewrites of
differing strategies and merges their results~\cite{gao2024dmqr}, and RaFe
trains a rewriter from reranker feedback rather than from generation-side
signals~\cite{ma2023rafe}.

Framed this way, rewriting reads as a monotone improvement. The classical
literature says otherwise. Mitra et al.\ named the failure ``query drift'':
expansion terms drawn from a first-pass set containing non-relevant documents
pull the expanded query toward those documents, so average gains coexist with
substantial per-query losses~\cite{mitra1998expansion}. Amati et al.\ showed the
losses concentrate on queries that are already difficult, and applied expansion
selectively on the basis of predicted difficulty~\cite{amati2004drift}.
Collins-Thompson treated expansion as a risk--reward problem, optimizing term
weights under constraints that bound the downside rather than maximize the
mean~\cite{collinsthompson2009risk}. The LLM-era analogue is Weller et al., who
evaluate generative query and document expansion across methods, retrievers, and
datasets and report ``a strong negative correlation between retriever
performance and gains from expansion: expansion improves scores for weaker
models, but generally harms stronger models,'' attributing the harm to noise
that expansions add alongside their recall
benefit~\cite{weller2024expansions}.

This body of work is the closest prior art to the finding reported here, and the
finding should be read against it. That expansion can degrade retrieval is
established, and the standard explanation --- that it shifts the query's term
distribution away from the relevant documents --- is consistent with what we
observe. Three things differ. First, the drift results above are measured over a
fixed collection with a fixed index, where every document is scored against
whatever query is issued, so drift changes the \emph{ranking} of a candidate set
that still contains the relevant documents. Retrieval in the present setting is a
set of calls to external search endpoints that each return a small,
provider-determined pool, so the phrasing determines which documents enter the
pipeline at all; drift becomes a fetch-time loss rather than a ranking-time loss,
and no downstream component can repair it. Second, the methods above augment
rather than replace: query2doc concatenates, DMQR-RAG merges multiple rewrites,
pseudo-relevance feedback re-weights an existing query vector. The failure
reported here arises from substitution --- an agent whose output replaces the
user's wording at every retrieval endpoint --- which is a property of how a
rewriting agent is wired into a pipeline, not of any rewriting method. Third, we
localize the loss by counting documents rather than by observing that a score
fell: 22 of the 23 relevant documents present in the baseline's retrieved set
were never fetched by the multi-agent pipeline, and three successively stronger
rerankers do not recover them. Prior work establishes that expansion can drift.
It does not establish where in an agentic pipeline the loss lands, nor that the
loss survives arbitrarily good ranking.

\subsection{Rewriting for Recall, Reranking for Precision}

Retrieval systems have long been staged: an inexpensive first pass over the
collection produces a candidate pool sized for recall, and an expensive second
pass reorders that pool for precision. BM25 and the probabilistic relevance
framework remain the reference first stage~\cite{robertson2009bm25}. The second
stage is where pretrained transformers were adopted first and most decisively.
monoBERT scores each query--document pair with a
cross-encoder~\cite{nogueira2019monobert}; monoT5 recasts the same task as
sequence-to-sequence relevance generation~\cite{nogueira2020monot5}; ColBERT
trades some cross-encoder accuracy for late interaction over precomputed term
embeddings~\cite{khattab2020colbert}; and RankGPT prompts a general-purpose
language model to reorder a candidate list
directly~\cite{sun2023rankgpt}. Lin et al.\ survey the pattern and its
variants~\cite{lin2021pretrained}.

The pattern carries an invariant that is easy to state and easy to lose: the
first stage bounds recall, and the second stage can only permute what the first
stage returned. Systems built on it therefore pair recall-oriented query handling
--- expansion, multiple rewrites, hybrid sparse--dense retrieval --- with a
precision-oriented reranker, and combine evidence from several query formulations
\emph{before} ranking. Reciprocal rank fusion is the standard cheap
combiner~\cite{cormack2009rrf}; DMQR-RAG's merge over multiple rewrites is the
same move inside a RAG pipeline~\cite{gao2024dmqr}.

The remedy reported in Section~\ref{sec:results} is an instance of this pattern,
not a new technique: issue both the original and the rewritten phrasing, union
the returned pools, then rank. Its interest is that the architecture it repairs
had adopted the recall half of the pattern (a rewriting agent) and the precision
half (an evaluator agent, and rerankers in the ablations) while omitting the
union that makes the two compatible, and that the omission was invisible in the
aggregate score the system was tuned against. The ranking ablation measures the
invariant directly: MRR rises from 0.250 to 0.625 across three rerankers, because
ranking within the fetched pool does improve, while nDCG@3 stays near 0.19,
because the pool does not contain the documents the metric is asking for.

\subsection{Multi-Agent RAG and Reproducibility Infrastructure}

The self-correcting designs reviewed above --- Self-RAG's reflection
tokens~\cite{asai2024selfrag}, CRAG's lightweight retrieval
evaluator~\cite{yan2024crag}, Adaptive-RAG's complexity-based
routing~\cite{jeong2024adaptiverag}, and FLARE's confidence-triggered retrieval
during generation~\cite{jiang2023flare} --- decide \emph{whether} and \emph{how
often} to retrieve, and how to filter what returns. None of them, as published,
checks whether the text handed to the retriever preserves the user's own
phrasing.

Two toolkits now make much of this work runnable. FlashRAG is an MIT-licensed
Python toolkit that reimplements RAG algorithms under one interface --- 16
methods over 38 datasets as published~\cite{jin2025flashrag}, 23 algorithms over
36 in the current release, including Self-RAG, Adaptive-RAG, and FLARE
pipelines~\cite{flashrag2025repo}. RAGLAB reproduces six algorithms across ten
benchmark datasets~\cite{zhang2024raglab} and publishes fine-tuned
\texttt{selfrag\_llama3-8B} checkpoints with VLLM-based and quantized generation
configurations~\cite{raglab2024repo}.

We did not run a head-to-head comparison against these baselines, for a reason
that is a property of our setting rather than a criticism of the toolkits: each
assumes a static corpus embedded and indexed in advance. Self-RAG's released
retrieval component is Contriever~\cite{izacard2022contriever} over the DPR
Wikipedia passage split~\cite{karpukhin2020dpr} with precomputed embeddings; its
authors note that retrieval over full English Wikipedia requires large memory
and multiple GPUs, and report testing their on-demand retrieval script on a
configuration with 100\,GB of RAM, while observing that less should
suffice~\cite{selfrag2023repo}. RAGLAB's ColBERT retrieval server is documented
as requiring at least 60\,GB of RAM~\cite{raglab2024repo}. CRAG's retrieval
evaluator is a fine-tuned T5-large model, and its confidence thresholds are set
separately for each evaluation dataset --- $(0.59, -0.99)$ for PopQA and
$(0.95, -0.91)$ for biography generation, among others~\cite{yan2024crag} --- so
transferring the method to a new domain requires re-tuning against labelled data
for that domain. The system studied here retrieves from live release, advisory,
and community APIs whose contents change between runs and for which no pre-built
index exists: there is no static split to embed, and per-dataset thresholds have
no analogue. The comparison we can make is against a single-agent baseline
issuing the same calls to the same endpoints in the same time window, which
isolates the architectural variable at the cost of not situating either system
against published leaderboard numbers. We record this as a limitation rather than
resolving it.

\subsection{Evaluation Methodology for Retrieval-Augmented Systems}

Retrieval evaluation inherits from TREC the use of pooled relevance judgments:
candidates are collected from the runs of all participating systems and judged
once, so no system is scored against a judgment set built around its own output.
Pooling is nonetheless incomplete. Zobel found that TREC measured relative system
performance reliably but overestimated recall, because many relevant documents
had likely never been found~\cite{zobel1998reliable}, and Buckley and Voorhees
showed that standard measures are not robust to substantially incomplete
judgments~\cite{buckley2004incomplete}. Both concerns apply to a comparison
between pipelines that fetch different document sets from live endpoints, and
they are why the judgments used here are pooled across both systems and all
ablation arms.

Benchmarks supply the other half of the practice. BEIR measures zero-shot
transfer across heterogeneous retrieval tasks~\cite{thakur2021beir}; KILT fixes a
single Wikipedia snapshot and requires provenance~\cite{petroni2021kilt}; the
CRAG benchmark spans five domains and eight question categories, with mock APIs
simulating web and knowledge-graph search~\cite{yang2024}; and FreshQA targets
questions whose answers change over time, the closest published analogue to
software update questions~\cite{vu2024freshllms}. Reference-free scorers such as
RAGAS~\cite{es2024ragas} and ARES~\cite{saadfalcon2024ares} estimate
faithfulness and relevance without gold answers.

Two aspects of that practice bear on the metric inversion reported here. The
first is judge independence. LLM judges agree with human preferences at rates
comparable to human--human agreement, but the same study documents position,
verbosity, and self-enhancement biases~\cite{zheng2023judge}, and evaluators
recognize and favor text they themselves generated~\cite{panickssery2024self}. A
judge that is also the pipeline's generator, or that is prompted with the
pipeline's own notion of relevance, is not independently measuring it. The second
is abstention. The CRAG benchmark scores an accurate answer $1$, a missing answer
$0$, and an incorrect answer $-1$, explicitly preferring an admission of
ignorance to a confident
error~\cite{yang2024}; unanswerable-question benchmarks draw the same
distinction on the reading-comprehension side~\cite{rajpurkar2018squad2}. A
keyword-overlap score registers neither distinction: it rewards the presence of
question terms in the retrieved text whether or not that text answers the
question, and assigns no cost to a fluent answer assembled from documents that do
not support one.

The general caution is not new. Armstrong et al.\ analysed a decade of reported
ad-hoc retrieval improvements and found little evidence of cumulative progress,
because baselines were generally weak --- often below the median original TREC
system --- and questioned the value of a statistically significant gain over such
a baseline~\cite{armstrong2009improvements}; Lipton and Steinhardt catalogue the
reporting patterns that let such results stand~\cite{lipton2019trends}. A
retrieval-quality score defined by the authors of the system it evaluates is the
limiting case, having neither external calibration nor a cross-system pool. What
this paper adds is a worked instance in which such a score and a standard one
disagree by roughly a factor of five in opposite directions over the same runs,
with the document-level audit that establishes which was right.

\subsection{Negative and Mixed Results}

Reporting that an architecture does not deliver what an earlier measurement
credited it with has precedent in software engineering and adjacent fields.
Menzies et al.\ published negative results for software effort estimation in
\emph{Empirical Software Engineering} under that heading, rather than reframing
the outcome~\cite{menzies2017negative}. Hellendoorn and Devanbu found that a
carefully engineered n-gram model outperformed the deep models then displacing
it for source code modeling, and identified properties of code that the neural
models handled poorly~\cite{hellendoorn2017dnn}. Ferrari Dacrema et al.\
examined 18 neural recommendation algorithms from top-level conferences, could
reproduce 7 with reasonable effort, and found that 6 of those 7 were often
outperformed by comparably simple heuristic
baselines~\cite{dacrema2019progress}. Liu et al.\ reviewed 147 deep learning
studies across 20 software engineering and 20 artificial intelligence venues and
found reproducibility and replicability routinely treated as minor threats or
deferred~\cite{liu2021reproducibility}. What these share with the present paper
is not the sign of the result but its structure: a measurement that had been
accepted is re-run under a protocol chosen independently of the system, the
ordering changes, and the mechanism behind the change is identified precisely
enough to transfer to other systems.
